\subsection{Site-dependent chains}

The source states its Lemma~5 for site-dependent tensors: five sites
carrying tensors $A_1, \ldots, A_5$ whose two-site windows are injective.
The entries above specialize to one repeated tensor; the remaining entries
of this section formalize the site-dependent setting itself.  A closed
chain carries one tensor per site, windows are arcs of consecutive sites,
and the closed-chain corollary delivers one gauge per bond.  Throughout,
all sites share one physical dimension $d$ and all bonds one bond
dimension $D$; the source poses no restriction on the local dimensions of
the site-dependent tensors, and the uniform-dimension restriction is
recorded in the route note
(\path{docs/paper-gaps/peps_normal_ft_section3_route.tex}).

\begin{definition}[Arc products and window injectivity of a site-dependent
        chain]%
    \label{def:chain_arc_window_injective}
    \lean{MPSChainTensor.arcEval}
    \lean{MPSChainTensor.IsWindowInjective}
    \leanok
    \uses{def:non_ti_chain}
    Let a closed chain on $n \ge 1$ sites carry one tensor of $D \times D$
    matrices over physical dimension $d$ per site, $A_0, \dots, A_{n-1}$.
    The \emph{arc product} of a word $x = x_1 \cdots x_\ell$ read from
    site $s$ is
    \[
        A_{(s)}^{x} = A_{s}^{x_1}\, A_{s+1}^{x_2} \cdots
        A_{s+\ell-1}^{x_\ell},
    \]
    site indices cyclic.  The chain is \emph{window injective at length
        $L$} when, for every site $s$, the arc products of the length-$L$
    words read from $s$ span the full matrix algebra.  This is the
    site-dependent setting of the closing section of
    \cite{Molnar2018NormalPEPS} (``Normal MPS: alternative proof''): an
    MPS on five sites in which the blocking of any two consecutive tensors
    is injective, for a general window length $L$.
\end{definition}

\begin{lemma}[Injectivity is window injectivity at length one]%
    \label{lem:isWindowInjective_one_of_isInjective}
    \lean{MPSChainTensor.isWindowInjective_one_of_isInjective}
    \leanok
    \uses{def:chain_arc_window_injective, def:injective}
    If every tensor in a closed chain is injective, then the chain is
    window injective at length \(1\).  Indeed, the length-one arc products
    at site \(s\) are precisely the matrices of the tensor at site \(s\).
\end{lemma}
\begin{proof}\leanok
    The family of length-one words is identified with the physical alphabet.
    Under this identification the length-one arc product is exactly the
    corresponding local matrix, so the two spanning conditions coincide.
\end{proof}

\begin{lemma}[Arc spanning from window injectivity]%
    \label{lem:isWindowInjective_arc_span}
    \lean{MPSChainTensor.IsWindowInjective.arc_span}
    \leanok
    \uses{def:chain_arc_window_injective}
    Let a site-dependent chain on $n \ge 1$ sites be window injective at
    length $L$ with $0 < L$.  Then for every $m \ge L$ and every site $s$,
    the arc products of the length-$m$ words read from $s$ span the full
    matrix algebra.  This is the site-dependent form of the remark that
    any region of at least the blocking size is injective
    (\cite{Molnar2018NormalPEPS}, closing section,
    cf.\ Lemma~\ref{lem:isNBlkInjective_of_le}).
\end{lemma}
\begin{proof}\leanok
    \uses{def:chain_arc_window_injective}
    Write the identity as a combination of window products at the starting
    site and peel the leading letter of each: any matrix becomes a
    combination of terms $A_s^{i}\,(W P)$ with $W$ an arc product of
    length $L - 1$ read from $s + 1$.  Each factor $W P$ is a combination
    of length-$m$ arc products read from $s + 1$ by induction, so each
    term is a combination of length-$(m+1)$ arc products read from $s$.
\end{proof}

\begin{lemma}[Bond-operator extraction from overlapping windows,
        site-dependent form]%
    \label{lem:chain_overlapWindow_exists_bondOperator}
    \lean{MPSChainTensor.overlapWindow_exists_bondOperator}
    \lean{MPSChainTensor.eq_of_span_mul_left}
    \leanok
    \uses{def:chain_arc_window_injective}
    Let $0 < L$ with $2L + 1 \le n$, let the closed chain on $n$ sites
    carry one tensor of $D \times D$ matrices per site, window injective
    at length $L$, and fix a site $r$.  Let $C_0, \ldots, C_L$ be deformed
    window tensors around the bond $(r+L-1,\, r+L)$: each $C_j$ assigns a
    matrix to every word of length $L$, read as replacing the blocked
    tensor of the $L$-site window starting at site $r + j$.  Suppose the
    deformed states agree for consecutive window positions: for every
    $j < L$, every prefix $p$ of $j$ sites read from $r$, every assignment
    $u$ on the $L + 1$ sites covered by the two windows, and every suffix
    $s$ on the remaining sites,
    \[
        \tr\bigl(A_{(r)}^{p}\, C_j(u_1 \cdots u_L)\,
        A_{r+j+L}^{u_{L+1}}\, A_{(r+j+L+1)}^{s}\bigr)
        = \tr\bigl(A_{(r)}^{p}\, A_{r+j}^{u_1}\,
        C_{j+1}(u_2 \cdots u_{L+1})\, A_{(r+j+L+1)}^{s}\bigr).
    \]
    Then the deformation is a virtual operation on the bond: there is a
    matrix $X$ with
    \[
        C_0(a) = A_{(r)}^{a}\, X
        \quad\text{and}\quad
        C_L(b) = X\, A_{(r+L)}^{b}
    \]
    for all length-$L$ words $a$, $b$, and $X$ is unique.  This is
    Lemma~5 of \cite{Molnar2018NormalPEPS} in the site-dependent setting
    in which the source states it, with the deformed window tensors
    arising as physical operators contracted with the blocked windows
    they act on; Lemma~\ref{lem:overlapWindow_exists_bondOperator} is its
    constant specialization.
\end{lemma}
\begin{proof}\leanok
    \uses{lem:isWindowInjective_arc_span, def:chain_arc_window_injective}
    Comparing the windows at $j$ and $j + 1$ through the complementary
    arc of $n - L - 1 \ge L$ sites, whose arc products span by
    Lemma~\ref{lem:isWindowInjective_arc_span}, strips the state equality
    to the letter level:
    $C_j(u_1 \cdots u_L)\,A_{r+j+L}^{u_{L+1}} =
        A_{r+j}^{u_1}\,C_{j+1}(u_2 \cdots u_{L+1})$ on every window of
    $L + 1$ sites.  Chaining the $L$ relations along the $2L$ sites
    flanking the bond turns the leftmost window into the rightmost one:
    $C_0(a)\,A_{(r+L)}^{b} = A_{(r)}^{a}\,C_L(b)$ for all length-$L$
    words.  Writing the identity as a combination
    $\sum_b \alpha_b A_{(r+L)}^{b} = I$, the matrix
    $X := \sum_b \alpha_b\,C_L(b)$ satisfies $A_{(r)}^{a} X = C_0(a)$;
    the mirrored combination of the $C_0$-family collapses onto the same
    matrix, giving $X\,A_{(r+L)}^{b} = C_L(b)$, and stripping the
    spanning factors makes $X$ unique.
\end{proof}

\begin{lemma}[Arc transport between same-state chains]%
    \label{lem:exists_arcTransport}
    \lean{MPSChainTensor.exists_arcTransport}
    \leanok
    \uses{def:chain_arc_window_injective, def:non_ti_chain}
    Let two site-dependent chains $A$, $B$ on $n$ sites generate the same
    closed-chain state, let $k + q = n$, and suppose the $A$-arc products
    of length $k$ read from a site $s$ and of length $q$ read from
    $s + k$ each span the full matrix algebra.  Then there is a linear
    map $\Lambda$ of the matrix algebra with
    $\Lambda(A_{(s)}^{x}) = B_{(s)}^{x}$ for every word $x$ of length
    $k$.
\end{lemma}
\begin{proof}\leanok
    \uses{lem:exists_mpvTransport, def:chain_arc_window_injective}
    As for one repeated tensor (Lemma~\ref{lem:exists_mpvTransport}):
    pair the matrix algebra against the complementary arc products.  The
    pairing is injective on the $A$-side by spanning and trace
    nondegeneracy, and the two pairings match on the spanning arc family
    because both compute closed-chain coefficients of total length $n$,
    equal after rotating the trace of the closed chain to the starting
    site $s$.
\end{proof}

\begin{lemma}[The insertion correspondence is an algebra homomorphism,
        site-dependent form]%
    \label{lem:chain_exists_insertionHom}
    \lean{MPSChainTensor.exists_insertionHom}
    \lean{MPSChainTensor.insertionHom_unique}
    \leanok
    \uses{def:chain_arc_window_injective, def:non_ti_chain}
    Let $0 < L$ with $2L + 1 \le n$, let $A$ and $B$ be site-dependent
    chains on $n$ sites, each window injective at length $L$, generating
    the same closed-chain state, and fix a site $r$.  Then there is a
    linear map $\Phi$ of the matrix algebra, unital and multiplicative,
    such that for every matrix $X$ and every word $w$ of length $n$,
    \[
        \tr\bigl(X\, A_{(r+L)}^{w}\bigr)
        = \tr\bigl(\Phi(X)\, B_{(r+L)}^{w}\bigr),
    \]
    and $\Phi$ is the unique linear map with this pairing property.  This
    is the closing clause of Lemma~5 of \cite{Molnar2018NormalPEPS} — the
    maps $O_1 \mapsto X$ and $O_3^T \mapsto X$ are uniquely defined and
    are algebra homomorphisms — in the site-dependent setting of the
    source, with the algebra structure carried by the insertions: $\Phi$
    sends an insertion on the bond $(r+L-1,\, r+L)$ of the first chain to
    the bond operator extracted on the second.
\end{lemma}
\begin{proof}\leanok
    \uses{lem:chain_overlapWindow_exists_bondOperator,
        lem:exists_arcTransport, lem:isWindowInjective_arc_span}
    A virtual insertion $X$ on the bond is realized on each of the
    $L + 1$ overlapping windows around it by deformed window tensors
    placing $X$ after the first $L - j$ letters of the window starting at
    site $r + j$.  Each deformation is transported to the second chain by
    the arc transport at its own starting site
    (Lemma~\ref{lem:exists_arcTransport}); the same-state hypothesis
    makes the transported deformations generate one common state, and the
    bond-operator extraction
    (Lemma~\ref{lem:chain_overlapWindow_exists_bondOperator}) produces
    $\Phi(X)$.  Uniqueness of the bond operator makes $\Phi$ linear,
    unital and multiplicative; the pairing against all closed-chain words
    follows by expanding the insertion over the window products, and pins
    $\Phi$ through the spanning arcs of length $n$
    (Lemma~\ref{lem:isWindowInjective_arc_span}) and nondegeneracy of the
    trace pairing.
\end{proof}

\begin{lemma}[Conjugate arc products at every bond, site-dependent form]%
    \label{lem:chain_exists_conjugation_of_sameState}
    \lean{MPSChainTensor.exists_conjugation_of_sameState}
    \leanok
    \uses{def:chain_arc_window_injective, def:non_ti_chain}
    Let $0 < L$ with $2L + 1 \le n$, and let $A$ and $B$ be
    site-dependent chains on $n$ sites with a common positive bond dimension,
    each window injective at length
    $L$, generating the same closed-chain state.  Then for every site $p$
    there is an invertible matrix $Z$ with
    \[
        B_{(p)}^{w} = Z\, A_{(p)}^{w}\, Z^{-1}
    \]
    for every word $w$ of length $n$.  The statement for one repeated
    tensor is Lemma~\ref{lem:exists_conjugation_of_mpv_eq}.
\end{lemma}
\begin{proof}\leanok
    \uses{lem:chain_exists_insertionHom, thm:skolem_noether}
    The insertion correspondence at the bond
    (Lemma~\ref{lem:chain_exists_insertionHom}) is unital and nonzero,
    hence an automorphism of the full matrix algebra, and by
    Skolem--Noether (Theorem~\ref{thm:skolem_noether}) it is conjugation
    by an invertible $Z$; pairing insertions against all closed-chain
    words read from $p$ and using nondegeneracy of the trace pairing
    transfers the conjugation to the arc products.
\end{proof}

\begin{corollary}[Fundamental Theorem for site-dependent normal closed
        chains at $n \ge 2L + 1$]%
    \label{cor:fundamentalTheorem_normalMPSChain_of_overlap}
    \lean{TNLean.PEPS.fundamentalTheorem_normalMPSChain_of_overlap}
    \leanok
    \uses{def:chain_arc_window_injective, def:non_ti_chain}
    Let $0 < L$ with $2L + 1 \le n$, and let $A$ and $B$ be
    site-dependent chains on $n$ sites with a common positive bond dimension,
    each window injective at length
    $L$, generating the same closed-chain state.  Then the chains are
    gauge equivalent: there are invertible matrices $Z_v$, one per bond
    of the closed chain, with
    \[
        B_v^i = Z_v\, A_v^i\, Z_{v+1}^{-1}
    \]
    for every site $v$ and every $i$, indices cyclic.  This is the
    closed-chain corollary of the closing section of
    \cite{Molnar2018NormalPEPS} assembled from its site-dependent
    Lemma~5 at every bond; the source displays the corollary for
    translation-invariant tensors
    (Corollary~\ref{cor:fundamentalTheorem_normalMPS_translationInvariant_of_overlap}),
    with the site-dependent Lemma~5 as its engine.
\end{corollary}
\begin{proof}\leanok
    \uses{lem:chain_exists_conjugation_of_sameState,
        lem:chain_overlapWindow_exists_bondOperator,
        lem:isWindowInjective_arc_span}
    The per-bond conjugations
    (Lemma~\ref{lem:chain_exists_conjugation_of_sameState}) give window
    covariance with nonzero scalars: on every arc of $m$ sites with
    $L \le m$ and $m + L \le n$, the conjugation at the near end
    intertwines the arc products of the two chains across the far bond,
    the bond-operator extraction
    (Lemma~\ref{lem:chain_overlapWindow_exists_bondOperator}) produces
    the connecting matrix, and the conjugation at the far end pins it to
    a nonzero scalar multiple of the gauge there through the centralizer
    of the full matrix algebra, so
    $B_{(p)}^{u} = c_{p,m} \cdot Z_p\, A_{(p)}^{u}\, Z_{p+m}^{-1}$ on all
    length-$m$ words $u$.  Comparing the covariance on windows of lengths
    $L + 1$ and $L$ through the spanning window products leaves the
    letterwise relation $B_v^i = \mu_v\, Z_v\, A_v^i\, Z_{v+1}^{-1}$ with
    nonzero scalars $\mu_v$; iterating it once around the closed chain
    against a nonzero arc product forces $\prod_v \mu_v = 1$, and
    dressing the gauges with the partial products of the scalars absorbs
    them, including across the seam.
\end{proof}

\begin{theorem}[Cyclic gauges preserve closed-chain states]%
    \label{thm:mps_chain_gaugeEquiv_sameState}
    \lean{MPSChainTensor.arcEval_eq_gauge_conj}
    \lean{MPSChainTensor.GaugeEquiv.sameState}
    \leanok
    \uses{def:chain_arc_window_injective, def:non_ti_chain}
    Let \(A_0,\ldots,A_{n-1}\) and \(B_0,\ldots,B_{n-1}\) be two
    site-dependent closed chains with a cyclic gauge relation
    \[
        B_v^i = Z_v\, A_v^i\, Z_{v+1}^{-1}
    \]
    at every site.  Then they generate the same closed-chain state.  More
    generally, the product along any arc is conjugated by the two gauges on
    the boundary bonds of that arc; for a full turn around the chain the
    two boundary bonds coincide, and cyclicity of trace removes the
    remaining conjugation.
\end{theorem}
\begin{proof}\leanok
    Iterate the sitewise gauge relation along the arc.  Consecutive
    boundary matrices cancel, leaving only the gauges at the two ends.
    For a word of length \(n\) read from the first site, the two ends are
    the same bond, so the closed trace is invariant under the resulting
    conjugation.
\end{proof}

\begin{corollary}[Fundamental Theorem for injective closed MPS chains]%
    \label{cor:fundamentalTheorem_injectiveMPSChain_of_sameState}
    \lean{TNLean.PEPS.fundamentalTheorem_injectiveMPSChain_of_sameState}
    \leanok
    \uses{cor:fundamentalTheorem_normalMPSChain_of_overlap,
        lem:isWindowInjective_one_of_isInjective}
    Let \(A_0,\ldots,A_{n-1}\) and \(B_0,\ldots,B_{n-1}\) be two
    site-dependent closed chains on \(n \ge 3\) sites, with common physical
    dimension and common positive bond dimension.  Suppose every local tensor in
    both chains is injective and the two closed-chain states are equal.
    Then there are invertible matrices \(Z_0,\ldots,Z_{n-1}\), one per
    bond of the closed chain, such that
    \[
        B_v^i = Z_v\, A_v^i\, Z_{v+1}^{-1}
    \]
    for every site \(v\) and every physical index \(i\), with indices read
    cyclically.  This is the \(L=1\), uniform-dimension instance of
    \cite[Theorem~\textup{thm:inj\_MPS}, lines~688--725]{Molnar2018NormalPEPS}.
\end{corollary}
\begin{proof}\leanok
    \uses{cor:fundamentalTheorem_normalMPSChain_of_overlap,
        lem:isWindowInjective_one_of_isInjective}
    Apply Corollary~\ref{cor:fundamentalTheorem_normalMPSChain_of_overlap}
    with \(L=1\).  The size condition becomes \(n \ge 3\), and
    Lemma~\ref{lem:isWindowInjective_one_of_isInjective} converts the
    sitewise injectivity hypotheses into the required window-injectivity
    hypotheses.
\end{proof}

\begin{definition}[Cyclic shift of a closed MPS chain]%
    \label{def:mps_chain_cyclic_shift}
    \lean{MPSChainTensor.cyclicShift}
    \lean{MPSChainTensor.IsCyclicShiftInvariantState}
    \leanok
    For a closed chain \(A_0,\ldots,A_{n-1}\), its cyclic shift is the
    chain whose \(v\)-th tensor is \(A_{v+1}\), with indices read modulo
    \(n\).  The generated state is cyclic-shift invariant when the
    closed-chain state of \(A_0,\ldots,A_{n-1}\) agrees with that of
    \(A_1,\ldots,A_{n-1},A_0\).
\end{definition}

\begin{corollary}[Gauge comparison with the cyclic shift]%
    \label{cor:fundamentalTheorem_injectiveMPSChain_cyclicShift}
    \lean{TNLean.PEPS.fundamentalTheorem_injectiveMPSChain_cyclicShift}
    \leanok
    \uses{cor:fundamentalTheorem_injectiveMPSChain_of_sameState,
        def:mps_chain_cyclic_shift}
    Let \(A_0,\ldots,A_{n-1}\) be an injective site-dependent closed chain
    with positive bond dimension on \(n\ge 3\) sites.  If its generated state
    is invariant under the
    cyclic shift \(A_v\mapsto A_{v+1}\), then there are invertible matrices
    \(Z_0,\ldots,Z_{n-1}\), one per bond, such that
    \[
        A_{v+1}^i = Z_v\, A_v^i\, Z_{v+1}^{-1}
    \]
    for every site \(v\) and physical index \(i\).  This is the first
    comparison step in the Applications-section translation-invariant
    description corollary of \cite[lines~1807--1824]{Molnar2018NormalPEPS}.
\end{corollary}
\begin{proof}\leanok
    Apply Corollary~\ref{cor:fundamentalTheorem_injectiveMPSChain_of_sameState}
    to the chain and its cyclic shift.  Injectivity is preserved by the
    shift, and the state equality is exactly the cyclic-shift invariance
    hypothesis.
\end{proof}

\begin{corollary}[Telescoping against the first tensor]%
    \label{cor:exists_gauge_to_first_of_cyclicShift_gaugeEquiv}
    \lean{MPSChainTensor.exists_gauge_to_first_of_cyclicShift_gaugeEquiv}
    \leanok
    \uses{def:mps_chain_cyclic_shift}
    Suppose a closed chain \(A_0,\ldots,A_{n-1}\) is gauge equivalent to its
    cyclic shift.  Then for every site \(v\) there are invertible matrices
    \(L_v\) and \(R_v\) such that
    \[
        A_v^i = L_v\, A_0^i\, R_v^{-1}
    \]
    for every physical index \(i\).  This is the telescoping step in
    \cite[lines~1828--1862]{Molnar2018NormalPEPS}, where the products denoted
    \(L_i\) and \(R_i\) express all tensors through the first one.
\end{corollary}
\begin{proof}\leanok
    Write the cyclic-shift gauge relation from one site to the next and
    iterate it along the chain.  At the first site one takes
    \(L_0=R_0=I\); passing from \(v\) to \(v+1\) multiplies the left and
    right virtual factors by the adjacent gauges.
\end{proof}

\begin{corollary}[Gauge-to-first description from a cyclic-invariant state]%
    \label{cor:exists_gauge_to_first_of_cyclicShiftInvariantState}
    \lean{TNLean.PEPS.exists_gauge_to_first_of_cyclicShiftInvariantState}
    \leanok
    \uses{cor:fundamentalTheorem_injectiveMPSChain_cyclicShift,
        cor:exists_gauge_to_first_of_cyclicShift_gaugeEquiv}
    Let \(A_0,\ldots,A_{n-1}\) be an injective site-dependent closed chain
    with positive bond dimension on \(n\ge 3\) sites.  If its generated state
    is invariant under the
    cyclic translation of the sites, then for every site \(v\) there are
    invertible matrices \(L_v\) and \(R_v\) such that
    \[
        A_v^i = L_v\, A_0^i\, R_v^{-1}
    \]
    for every physical index \(i\).  This is the uniform-bond-dimension form
    of the \(L_v,R_v\) step in the Applications-section proof of
    \cite[lines~1803--1890]{Molnar2018NormalPEPS}.
\end{corollary}
\begin{proof}\leanok
    Apply the injective chain theorem to the chain and its cyclic shift, then
    telescope the resulting bond gauges around the chain.
\end{proof}

\begin{corollary}[Translation-invariant description of an injective closed MPS chain,
        uniform bond dimension]%
    \label{cor:exists_constant_injectiveMPS_of_cyclicShiftInvariantState}
    \lean{TNLean.PEPS.exists_constant_injectiveMPS_of_cyclicShiftInvariantState}
    \leanok
    \uses{cor:fundamentalTheorem_injectiveMPSChain_cyclicShift,
        def:mps_chain_cyclic_shift}
    Let \(A_0,\ldots,A_{n-1}\) be an injective site-dependent closed chain on
    \(n\ge 3\) sites, with one physical dimension \(d\) and one positive bond
    dimension \(D\).  If the closed-chain state is invariant under the cyclic shift
    \(A_v\mapsto A_{v+1}\), then there is an injective tensor \(B\), with bond
    dimension \(D\), such that for every physical configuration \(\sigma\),
    \[
        \tr\!\left(
            A_0^{\sigma_0}\cdots A_{n-1}^{\sigma_{n-1}}\right)
        =
        \tr\!\left(
            B^{\sigma_0}\cdots B^{\sigma_{n-1}}\right).
    \]

    This is the uniform-bond-dimension target corresponding to the
    Applications-section corollary of
    \cite[lines~1803--1890]{Molnar2018NormalPEPS}.  The source also concludes
    that all bond dimensions of the original site-dependent representation are
    equal; in the present uniform-dimension statement that clause has already
    been built into the setting.
\end{corollary}
\begin{proof}\leanok
    Apply Corollary~\ref{cor:fundamentalTheorem_injectiveMPSChain_cyclicShift}
    to compare the chain with its cyclic shift, giving one invertible gauge
    $Z_v$ per bond with $A_{v+1}^i=Z_v\,A_v^i\,Z_{v+1}^{-1}$.  Telescoping
    the gauges into the running products $P_m=Z_{m-1}\cdots Z_0$ expresses
    every tensor through the first, $A_m^i=P_m\,(A_0^i Z_0)\,P_{m+1}^{-1}$, so
    the closed trace collapses to
    $\tr(B^{\sigma_0}\cdots B^{\sigma_{n-1}} P_n^{-1})$ with the
    tentative repeated tensor $B=A_0 Z_0$.  Reading the gauge relation across
    the seam $A_{n+1}\equiv A_0$ shows that the loop product $P_n$ commutes
    with the full matrix algebra spanned by $B$, hence is a scalar
    $\lambda\ne 0$.  Since $\mathbb{C}$ is algebraically closed there is an
    $n$-th root $c^n=\lambda^{-1}$; the rescaled tensor $c\,B$ is still
    injective and its constant chain generates the same closed state.
\end{proof}

\begin{corollary}[Translation-invariant description of an injective 2D PEPS,
        uniform bond dimension]%
    \label{cor:exists_constant_injectivePEPS_of_translationInvariantState}
    \notready
    \uses{cor:exists_constant_injectiveMPS_of_cyclicShiftInvariantState,
        thm:fundamentalTheorem_PEPS, def:peps_torus_translation_invariant,
        def:peps_torus_uniform_bond_dim,
        thm:peps_stateCoeff_translationInvariant}
    Let \(A\) be an injective PEPS on the discrete torus with orientation-uniform
    bond dimensions, horizontal dimension \(D_h\) and vertical dimension \(D_v\),
    and positive bond dimensions.  If the generated state is invariant under all
    lattice translations \(\tau_{a,b}\), that is, for every translation and every
    physical configuration \(\sigma\),
    \[
        \Coeff_{A}(\sigma\circ\tau_{a,b})
        =
        \Coeff_{A}(\sigma),
    \]
    then there is an injective torus tensor \(B\), with the same orientation-uniform
    bond dimensions \(D_h\) and \(D_v\), whose constant (site-independent) family
    generates the same state: for every physical configuration \(\sigma\),
    \[
        \Coeff_{A}(\sigma)
        =
        \Coeff_{B}(\sigma).
    \]

    This is the uniform-bond-dimension target corresponding to the
    Applications-section 2D corollary of
    \cite[lines~1892--1894]{Molnar2018NormalPEPS}.  The source additionally
    concludes that all vertical bond dimensions agree and all horizontal bond
    dimensions agree; in the present orientation-uniform setting those two
    clauses have already been built into the bond-dimension type, so they carry
    no content here.
\end{corollary}
\begin{proof}
    The source states this corollary without proof, remarking only that the
    one-dimensional argument extends; the route below is that intended
    extension.  Write \(A^{(j,k)}\) for the site tensor at lattice position
    \((j,k)\), with horizontal virtual legs of dimension \(D_h\) and vertical
    virtual legs of dimension \(D_v\).  Comparing the state with its
    \(\tau_{1,0}\)-translate through the injective PEPS Fundamental Theorem
    (Theorem~\ref{thm:fundamentalTheorem_PEPS}) produces, on each horizontal
    edge, invertible \(D_h\times D_h\) gauges \(X^{(j,k)}_h\) with
    \[
        A^{(j+1,k)} = \bigl(X^{(j,k)}_h\bigr)^{-1}\, A^{(j,k)}\, X^{(j+1,k)}_h,
    \]
    acting on the two horizontal virtual legs, and comparing with the
    \(\tau_{0,1}\)-translate produces, on each vertical edge, invertible
    \(D_v\times D_v\) gauges \(X^{(j,k)}_v\) with
    \[
        A^{(j,k+1)} = \bigl(X^{(j,k)}_v\bigr)^{-1}\, A^{(j,k)}\, X^{(j,k+1)}_v.
    \]
    Telescoping the horizontal gauges along a row reduces each site to the
    leftmost one,
    \[
        A^{(j,k)} = \bigl(L^{(j,k)}_h\bigr)^{-1}\, A^{(0,k)}\, R^{(j,k)}_h,
        \qquad
        L^{(j,k)}_h = X^{(0,k)}_h\cdots X^{(j-1,k)}_h,
        \quad
        R^{(j,k)}_h = X^{(1,k)}_h\cdots X^{(j,k)}_h,
    \]
    and the periodic boundary condition closes the row product into a single
    residual virtual operation, the 2D analogue of the one-dimensional identity
    \(R_vL_{v+1}^{-1}=Z_0^{-1}\):
    \[
        R^{(j,k)}_h\,\bigl(L^{(j+1,k)}_h\bigr)^{-1} = \bigl(X^{(0,k)}_h\bigr)^{-1}.
    \]
    The same telescoping closes each column, yielding one residual vertical
    operation.  Absorbing the two residual gauges into the corner site gives a
    single repeated tensor
    \[
        B = \bigl(X^{(0,0)}_v\bigr)^{-1}\,\bigl(X^{(0,0)}_h\bigr)^{-1}\, A^{(0,0)},
    \]
    whose constant family generates the same state.  The one-dimensional
    telescoping closure is stated as
    Corollary~\ref{cor:exists_constant_injectiveMPS_of_cyclicShiftInvariantState}.
    What remains is the two-directional PEPS assembly.  If
    \[
        H_{j,k}=R^{(j,k)}_h\bigl(L^{(j+1,k)}_h\bigr)^{-1},
        \qquad
        V_{j,k}=R^{(j,k)}_v\bigl(L^{(j,k+1)}_v\bigr)^{-1},
    \]
    are the residual horizontal and vertical operations, the horizontal residual
    acts on \(\mathbb C^{D_h}\) and the vertical residual acts on
    \(\mathbb C^{D_v}\).  The missing plaquette flatness identity must therefore
    be stated after placing them on the corner space
    \(\mathbb C^{D_h}\otimes\mathbb C^{D_v}\).  Equivalently, after the row and
    column telescopings make \(H_{j,k}=H_k\) and \(V_{j,k}=V_j\), one must prove
    \[
        (H_k\otimes I_{D_v})(I_{D_h}\otimes V_{j+1})
        =
        \lambda_{j,k}\,(I_{D_h}\otimes V_j)(H_{k+1}\otimes I_{D_v})
    \]
    for nonzero scalars \(\lambda_{j,k}\).  The scalar phases must also close
    around the two periodic seams,
    \[
        \prod_{k\in\mathbb Z/N_v\mathbb Z}\lambda_{j,k}=1,
        \qquad
        \prod_{j\in\mathbb Z/N_h\mathbb Z}\lambda_{j,k}=1.
    \]
    These identities are then absorbed into this single repeated 2D tensor.
\end{proof}

\begin{theorem}[Uniqueness of the injective closed-chain gauge]%
    \label{thm:fundamentalTheorem_injectiveMPSChain_gauge_unique}
    \lean{TNLean.PEPS.fundamentalTheorem_injectiveMPSChain_gauge_unique}
    \leanok
    \uses{cor:fundamentalTheorem_injectiveMPSChain_of_sameState,
        def:chain_arc_window_injective}
    Under the positive-bond-dimension hypotheses of
    Corollary~\ref{cor:fundamentalTheorem_injectiveMPSChain_of_sameState},
    suppose \(Z_v\) and \(Z'_v\) are two families of invertible matrices
    satisfying
    \[
        B_v^i = Z_v\, A_v^i\, Z_{v+1}^{-1}
        \quad\text{and}\quad
        B_v^i = Z'_v\, A_v^i\, (Z'_{v+1})^{-1}
    \]
    for every \(v,i\).  Then there is a single nonzero scalar \(c\) such
    that \(Z'_v = c Z_v\) for every bond \(v\).  This is the uniqueness
    clause in \cite[line~724]{Molnar2018NormalPEPS}.
\end{theorem}
\begin{proof}\leanok
    For each bond put \(C_v = Z_v^{-1}Z'_v\).  Comparing the two gauge
    relations gives
    \[
        C_v A_v^i = A_v^i C_{v+1}
    \]
    for every local matrix.  Since the matrices of \(A_v\) span the full
    matrix algebra, evaluating this identity on the identity matrix gives
    \(C_v=C_{v+1}\); hence all \(C_v\) are one common matrix.  The same
    relation then says that this common matrix commutes with the full
    matrix algebra, so it is a scalar matrix.  Its invertibility makes the
    scalar nonzero.
\end{proof}

Two remarks place these statements against \cite{Molnar2018NormalPEPS}.  The
injective matrix product states of the opening section of
\cite{Molnar2018NormalPEPS} are the case $L = 1$, in which a single tensor
already spans the matrix algebra; the corollaries above contain them at
$n \ge 3$, with no separate argument needed.  The site-dependent results are
stated for one physical dimension $d$ and one bond dimension $D$ around the
chain, while \cite{Molnar2018NormalPEPS} places no restriction on the local
dimensions; this uniform-dimension restriction is the only narrowing.
% Paper-gap note: docs/paper-gaps/peps_normal_ft_section3_route.tex

\section{The vertex-injective case and the residual edge scalars}

\begin{theorem}[Fundamental Theorem for normal PEPS, vertex-injective case]%
    \label{thm:fundamentalTheorem_normalPEPS_of_vertexInjective}
    \lean{TNLean.PEPS.fundamentalTheorem_normalPEPS_of_vertexInjective}
    \leanok
    \uses{thm:fundamentalTheorem_PEPS, def:peps_normalPEPSBlockingHypotheses,
        def:GaugeEquivPEPS, def:IsVertexInjective}
    Suppose PEPS tensors $A$ and $B$ are vertex injective with positive bond
    dimensions on a connected graph, generate the same state, and $A$ satisfies
    the normal blocking hypotheses. Then their defining tensors are related by
    a local gauge.
    This is the vertex-injective specialization of the general normal-PEPS
    theorem: when each single tensor is already injective the gauge follows
    from the injective Fundamental Theorem, so the region blocking hypotheses
    are not used. The general region-injective theorem, in which only blocked
    regions are injective, remains open; the missing step is the
    region-level insertion-algebra isomorphism that produces the per-edge gauge
    after blocking.
\end{theorem}
\begin{proof}\leanok
    Vertex injectivity, positive bond dimensions, same state, and connectivity
    are exactly the hypotheses of the injective Fundamental Theorem, which
    supplies the local gauge.
\end{proof}

\begin{definition}[Oriented endpoint scalar]%
    \label{def:edgeScalarAt}
    \lean{TNLean.PEPS.edgeScalarAt}
    \leanok
    \uses{def:edgeGaugeAt}
    An edge-scalar family $c=(c_e)_{e\in E}$ determines the endpoint scalar
    $c_{v,e}$ by $c_{v,e}=c_e$ at the lower endpoint of $e$ and
    $c_{v,e}=c_e^{-1}$ at the upper endpoint.
\end{definition}

\begin{definition}[Vertex-balanced edge scalars]%
    \label{def:IsVertexBalanced}
    \lean{TNLean.PEPS.IsVertexBalanced}
    \leanok
    \uses{def:edgeScalarAt}
    An edge-scalar family is vertex-balanced when
    $\prod_{e\ni v} c_{v,e}=1$ for every vertex $v$.
\end{definition}
